\par
\subsection{Lorentzian transport and electromagnetic representation}

The electronic section provides a structure and a coordinate
before its application to an interaction. The Lorentzian realization
of cubic incidence determines, under the stated temporal convention,
the space of forms in which that interaction can be expressed.
To formulate it, an oriented local realization \((\mathcal M,g)\)
of dimension four is fixed, with tangent fibres identified with the model
\(Q_\square\). A connection \(\nabla\) with
\(\nabla g=0\), a potential one-form \(A\in\Omega^1(\mathcal M)\),
its curvature \(F=dA\), and a charge representation of the section
are also specified. These data are the arguments of the physical representation
that follows.

Hodge duality satisfies \(\star^2=-I\) on two-forms.
The electric--magnetic separation corresponds to a temporal choice
in the \(1+3\) decomposition; \(F\) and \(\star F\) preserve their
relationship within the same space of two-forms. Neutrality of a
charge representation means \(q=0\), whereas the section may
retain orientation, phase, and internal transport. The absence of
electric coupling in that representation has a different content
from the absence of phase information.

In this dimensional realization, let \(m_e>0\), let \(q\) be the charge,
let \(\gamma(s)\) be a parametrized trajectory, and let \(u=\dot\gamma\)
be its four-velocity. The Lorentz force is expressed by
\begin{equation}
 m_e\nabla_u u^\mu=qF^\mu{}_{\nu}u^\nu.
 \label{eq:lorentz-electron}
\end{equation}
The antisymmetry of \(F\) implies
\(F_{\mu\nu}u^\mu u^\nu=0\). Consequently,
\[
 \frac{d}{ds}g(u,u)=2g(\nabla_u u,u)=0.
\]
The normalization of the trajectory is preserved provided that \(\nabla\)
is metric and \eqref{eq:lorentz-electron} holds. The mass used
is the dimensional realization of the preceding electronic section;
the coupling parameter is identified in the adopted electromagnetic
convention. The concrete deduction of the potential and of a field
configuration constitutes an additional problem with its own
boundary conditions. The result presented here determines the domain, the
representation data, and the conservation implied by its equation.

TPK route inversion and the orientation of an electronic line
then offer a precise comparison: each has its involution,
transport, and representation. The work of Wheeler and Feynman
places the question historically; the equivalence of two realizations
is decided through their maps and compatibility conditions.
This arrangement keeps the electron's arithmetic origin connected
to the formulation of its interaction, making the function
of each stage explicit.
